% Part: computability
% Chapter: computability-theory
% Section: k-1

\documentclass[../../../include/open-logic-section]{subfiles}

\begin{document}

\olfileid{cmp}{thy}{k1}
\olsection{نمونه‌ای از کاهش‌پذیری}

کاربردی از \olref[ppr]{prop:reduce} را در نظر بگیریم.

\begin{prop}
\ollabel{prop:k1}
بگذارید
\[
K_1 = \Setabs{e}{\cfind{e}(0) \fdefined}.
\]
آنگاه $K_1$ محاسبه‌پذیرِ برشمردنی است، اما محاسبه‌پذیر نیست.
\end{prop}

\begin{proof}
چون $K_1=\Setabs{e}{\lexists[s][T(e,0,s)]}$، بنا بر
\olref[eqc]{thm:exists-char} مجموعهٔ $K_1$ محاسبه‌پذیرِ برشمردنی است.

برای نشان‌دادن اینکه $K_1$ محاسبه‌پذیر نیست، نشان می‌دهیم $K_0$ به
آن کاهش‌پذیر است.

\begin{explain}
این کار قدری دشوار است، زیرا با استفاده از $K_1$ تنها می‌توانیم
دربارهٔ محاسبه‌هایی پرسش کنیم که با ورودی مشخص~$0$ آغاز می‌شوند. فرض
کنید دوست باهوشی دارید که می‌تواند به پرسش‌های این‌چنینی پاسخ دهد
(چنین دوستانی را «اوراکل» می‌نامند). سپس فرض کنید کسی نزد شما می‌آید
و می‌پرسد آیا $\tuple{e,x}$ در $K_0$ است یا نه؛ یعنی آیا ماشین~$e$
روی ورودی~$x$ متوقف می‌شود. یک کار ممکن این است که ماشین دیگری
چون~$e_x$ بسازید که برای \emph{هر} ورودی، آن ورودی را نادیده بگیرد
و در عوض~$e$ را روی ورودی~$x$ اجرا کند. آنگاه روشن است که پرسش توقف
ماشین~$e$ روی ورودی~$x$ هم‌ارز پرسش توقف ماشین~$e_x$ روی ورودی~$0$
(یا هر ورودی دیگر) است. پس از دوستتان می‌پرسید آیا این ماشین تازه،
$e_x$، روی ورودی~$0$ متوقف می‌شود؛ پاسخ او به پرسش تغییریافته پاسخ
پرسش اصلی را به دست می‌دهد. این همان کاهش مطلوب $K_0$ به~$K_1$ است.
\end{explain}

با استفاده از تابع محاسبه‌پذیر جزئیِ جهانی، بگذارید $f$ تابع
سه‌موضعیِ زیر باشد:
\[
f(x,y,z) \simeq \cfind{x}(y).
\]
توجه کنید $f$ ورودی سومش را به‌کلی نادیده می‌گیرد. نمایه‌ای چون $e$
انتخاب کنید که $f=\cfind{e}[3]$؛ پس داریم
\[
\cfind{e}[3](x,y,z) \simeq \cfind{x}(y).
\]
بنا بر قضیهٔ $s$-$m$-$n$، تابعی چون $s(e,x,y)$ وجود دارد که برای
هر $z$،
\begin{align*}
\cfind{s(e,x,y)}(z) & \simeq \cfind{e}[3](x,y,z) \\
& \simeq \cfind{x}(y).
\end{align*}

\begin{explain}
برحسب استدلال غیررسمی بالا، $s(e,x,y)$ نمایه‌ای برای ماشینی است که
برای هر ورودی $z$، آن ورودی را نادیده می‌گیرد و $\cfind{x}(y)$ را
محاسبه می‌کند.
\end{explain}

به‌ویژه، داریم
\[
\cfind{s(e,x,y)}(0) \fdefined \quad \text{اگر و تنها اگر} \quad
\cfind{x}(y) \fdefined.
\]
به بیان دیگر، $\tuple{x,y}\in K_0$ اگر و تنها اگر
$s(e,x,y)\in K_1$. پس تابع $g$ که با
\[
g(w) = s(e,(w)_0,(w)_1)
\]
تعریف می‌شود، کاهشی از $K_0$ به~$K_1$ است.
\end{proof}

\end{document}
